\cotefile{58}
\scanpage{20}
\mspage{- 2 -}
\margin{n°358}

$V'_{(K)}$ dans $W$ et une seule telle que $\bar f'(v'\otimes 1)=f'(v')$ pour tout $v'$ dans $V'$.
Si $\lambda$ est un isomorphisme, l'application $K$-linéaire $\bar f'\mathbin{:}\lambda^{-1}$ est alors
l'unique application $K$-linéaire de $V$ dans $W$ prolongeant $f'$ ; autrement dit,
a) implique b). D'autre part, si b) est satisfaite, il existe une application $K$-linéaire $\mu$ de $V$ dans $V'_{(K)}$ telle que $\mu(v')=v'\otimes 1$ pour tout $v'$
dans $V'$ ; il est immédiat que $\lambda$ et $\mu$ sont des applications réciproques, et
par suite $\lambda$ est une bijection. On a donc prouvé que a) équivaut à b).

Par ailleurs, pour toute base $B'$ de $V'$ sur $K'$, l'ensemble $B$ des éléments
$v'\otimes 1$ de $V'_{(K)}$ (pour $v'$ parcourant $B'$) est une base de $V'_{\inserted{(K)}}$ sur $K$, et $\lambda$
transforme $B$ en $B'$. Pour que l'application $K$-linéaire $\lambda$ soit bijective, il
faut que l'image par $\lambda$ de toute base de $V'_{(K)}$ soit une base de $V$ sur $K$, et
il suffit qu'il en soit ainsi pour une seule base. Ceci prouve l'équivalence
de a), c) et d).

Comme toute partie de $V'$ libre sur $K'$ est contenue dans une base de $V'$
sur $K'$, c) implique e). Enfin, supposons e) satisfaite ; si $B'$ est une base
de $V'$ sur $K'$, c'est une partie de $V$ libre sur $K$, et comme $B'$ engendre $V'$
sur $K'$, donc $V$ sur $K$ d'après e), on voit que $B'$ est une base de $V$ sur $K$ ;
on a donc prouvé que e) implique c). C.Q.F.D.

DEFINITION 1.- \uline{Soient} $V$ \uline{un espace vectoriel à droite sur un corps} $K$ \uline{et} $K'$ \uline{un}
\uline{sous-corps de} $K$. \uline{On appelle} $K'$-\uline{structure sur} $V$ \uline{tout sous-espace} $K'$-\uline{vectoriel}
$V'$ \uline{de} $V$ \uline{qui vérifie les conditions équivalentes de la prop.1}. \uline{On dit que les}
\uline{éléments de} $V'$ \uline{sont les vecteurs rationnels sur} $K'$.

2. \uline{Rationalité pour un sous-espace.}

DEFINITION 2.- \uline{Soient} $V$ \uline{un espace vectoriel à droite sur le corps} $K$ \uline{et} $V'$
\uline{une} $K'$-\uline{structure sur} $V$. \uline{On dit qu'un sous-espace} $K$-\uline{vectoriel} $W$ \uline{de} $V$ \uline{est}
\uline{rationnel sur} $K'$ \uline{s'il est engendré sur} $K$ \uline{par} $W\cap V'$.

\margin{19}
